\PassOptionsToPackage{table}{xcolor}
\documentclass[10pt,aspectratio=169]{beamer}
\newcommand{\LectureNumber}{20}
\input{course-theme.tex}
\title{One-dimensional arrays}
\subtitle{CSC103 Programming Fundamentals / Lecture 20}
\author{Dr. Muhammad Farhan}
\institute{Associate Professor of Computer Science\\Department of Computer Science\\COMSATS University Islamabad\\Sahiwal Campus}
\date{Fall 2026}
\begin{document}
\begin{frame}\titlepage\end{frame}

\begin{frame}[fragile,t]{The problem for this lecture}
\begin{columns}[T,onlytextwidth]\begin{column}{0.60\textwidth}
Find the maximum in \{-4,-9,-2\}. State the policy for an empty array.
\end{column}\begin{column}{0.35\textwidth}\small
Require a nonempty array before accessing values[0].
\end{column}\end{columns}
\note{Introduce the target problem before explaining the tools. Revisit it in the guided solution later.\par Syllabus reading: Deitel Ch 7. CLO-2.}
\end{frame}

\begin{frame}[fragile,t]{Learning objectives}
\begin{columns}[T,onlytextwidth]\begin{column}{0.60\textwidth}
\begin{itemize}
\item Create arrays with fixed or runtime sizes
\item Traverse elements safely
\item Compute an aggregate and explain bounds errors
\end{itemize}
\end{column}\begin{column}{0.35\textwidth}\small
CLO-2

Array declaration and creation\par Indexes and length\par Array traversal\par Runtime size and aggregation
\end{column}\end{columns}
\note{Explain how each objective supports the opening problem. The final questions check these objectives.\par Syllabus reading: Deitel Ch 7. CLO-2.}
\end{frame}

\begin{frame}[fragile,t]{Array declaration and creation}
\begin{itemize}
\item An array groups indexed elements of one component type.
\item An array reference and the array object are different things.
\item The array length stays fixed after construction.
\end{itemize}
\vspace{1mm}\begin{block}{Example}
\begin{lstlisting}
int[] marks = {60, 75, 90};
int[] counts = new int[4];
\end{lstlisting}
\end{block}
\note{A new int array starts with zero-valued elements. A local array reference itself must still be assigned before use.\par Syllabus reading: Deitel Ch 7. CLO-2.}
\end{frame}

\begin{frame}[fragile,t]{Indexes and length}
\begin{itemize}
\item Indexes start at zero and end at length-1.
\item length is an array property, without parentheses.
\item An invalid index raises ArrayIndexOutOfBoundsException.
\end{itemize}
\vspace{1mm}\begin{block}{Example}
\begin{lstlisting}
marks[0]       // 60
marks.length   // 3
marks[3]       // invalid
\end{lstlisting}
\end{block}
\note{An empty array is valid and has no valid element index. Distinguish marks.length from a String length() call.\par Syllabus reading: Deitel Ch 7. CLO-2.}
\end{frame}

\begin{frame}[fragile,t]{Array traversal}
\begin{itemize}
\item An indexed loop supports reading and updating elements.
\item An enhanced for loop conveniently reads each value.
\item Keep the condition i \textless{} array.length.
\end{itemize}
\vspace{1mm}\begin{block}{Example}
\begin{lstlisting}
int total = 0;
for (int value : marks) {
  total += value;
}
\end{lstlisting}
\end{block}
\note{Assigning to the enhanced-for variable does not update primitive array elements. Use marks[i] for an element update.\par Syllabus reading: Deitel Ch 7. CLO-2.}
\end{frame}

\begin{frame}[fragile,t]{Runtime size and aggregation}
\begin{itemize}
\item Read or compute a size before new int[size].
\item Reject negative sizes and unreasonable requests.
\item A mean or maximum needs an explicit empty-array policy.
\end{itemize}
\vspace{1mm}\begin{block}{Example}
\begin{lstlisting}
int size = 5;
int[] data = new int[size];
// data has indexes 0..4
\end{lstlisting}
\end{block}
\note{new int[-1] throws NegativeArraySizeException. An extremely large allocation may fail. For maximum, initialise from data[0] only after ensuring nonempty input.\par Syllabus reading: Deitel Ch 7. CLO-2.}
\end{frame}

\begin{frame}[fragile,t]{The initial maximum must come from the domain}
\begin{itemize}
\item Zero is an incorrect default maximum when all values can be negative.
\item For a nonempty array, the first element supplies a valid candidate.
\item Each comparison then maintains the maximum of the processed prefix.
\end{itemize}
\vspace{1mm}\begin{block}{Example}
\begin{lstlisting}
Input: {-4,-9,-2}
Start at -4. Compare -9, then -2.
Final maximum: -2
\end{lstlisting}
\end{block}
\note{Explain the mechanism using the displayed example. Zero is an incorrect default maximum when all values can be negative. For a nonempty array, the first element supplies a valid candidate. Each comparison then maintains the maximum of the processed prefix.\par Syllabus reading: Deitel Ch 7. CLO-2.}
\end{frame}

\begin{frame}[fragile,t]{A traversal bound is part of correctness}
\begin{itemize}
\item For an array of length n, valid indexes satisfy 0 \textless{}= i \textless{} n.
\item Starting at 1 is appropriate after using index 0 as the initial candidate.
\item A length-zero input has no first element, so check it before initialisation.
\end{itemize}
\vspace{1mm}\begin{block}{Example}
\begin{lstlisting}
int maximum = values[0];
for (int i=1; i<values.length; i++) ...
\end{lstlisting}
\end{block}
\note{Explain the mechanism using the displayed example. For an array of length n, valid indexes satisfy 0 \textless{}= i \textless{} n. Starting at 1 is appropriate after using index 0 as the initial candidate. A length-zero input has no first element, so check it before initialisation.\par Syllabus reading: Deitel Ch 7. CLO-2.}
\end{frame}


\begin{frame}[fragile,t]{An array stores indexed elements}
\begin{center}\begin{tikzpicture}
\node[box,text width=13mm,minimum height=11mm] (n0) at (0.0,0.0) {12};
\node[font=\scriptsize,text=accent] at (0.0,.85) {index 0};
\node[box,text width=13mm,minimum height=11mm] (n1) at (1.8,0.0) {7};
\node[font=\scriptsize,text=accent] at (1.8,.85) {index 1};
\node[box,text width=13mm,minimum height=11mm] (n2) at (3.6,0.0) {19};
\node[font=\scriptsize,text=accent] at (3.6,.85) {index 2};
\node[box,text width=13mm,minimum height=11mm] (n3) at (5.4,0.0) {4};
\node[font=\scriptsize,text=accent] at (5.4,.85) {index 3};
\end{tikzpicture}\end{center}
\begin{block}{What to notice}A length-4 array has indexes 0 through 3. Index 4 is outside the array.\end{block}
\note{Trace each labelled connection. A length-4 array has indexes 0 through 3. Index 4 is outside the array.}
\end{frame}
\begin{frame}[fragile,t]{Key distinctions}
\small\renewcommand{\arraystretch}{1.5}
\rowcolors{2}{pale}{white}
\begin{tabularx}{\textwidth}{>{\raggedright\arraybackslash}X>{\raggedright\arraybackslash}X>{\raggedright\arraybackslash}X>{\raggedright\arraybackslash}X}
\rowcolor{navy}
\color{white}\bfseries Index & \color{white}\bfseries 0 & \color{white}\bfseries 1 & \color{white}\bfseries 2\\
Element & -4 & -9 & -2\\
Valid for length 3? & Yes & Yes & Yes\\
Candidate after processing & -4 & -4 & -2\\
\end{tabularx}
\vspace{3mm}\footnotesize These distinctions determine which operation or control structure fits the problem.
\note{\par Syllabus reading: Deitel Ch 7. CLO-2.}
\end{frame}

\begin{frame}[fragile,t]{First example: execution}
\begin{lstlisting}
int[] marks = {60, 75, 90};
int total = 0;
for (int i = 0; i < marks.length; i++) {
  total += marks[i];
}
System.out.println((double) total / marks.length);
\end{lstlisting}
\note{This is the main body from Java\_Examples/Lecture20.java. The source file includes the complete class. Predict the result before the next slide.\par Syllabus reading: Deitel Ch 7. CLO-2.}
\end{frame}

\begin{frame}[fragile,t]{First example: result}
\begin{columns}[T,onlytextwidth]\begin{column}{0.60\textwidth}
75.0\par The running sums are 60, 135 and 225.\par The loop uses indexes 0, 1 and 2.
\end{column}\begin{column}{0.35\textwidth}\small
Array traversal

Runtime size and aggregation
\end{column}\end{columns}
\note{Expected output and interpretation: 75.0
The running sums are 60, 135 and 225.
The loop uses indexes 0, 1 and 2.\par Syllabus reading: Deitel Ch 7. CLO-2.}
\end{frame}

\begin{frame}[fragile,t]{Guided practice}
\begin{columns}[T,onlytextwidth]\begin{column}{0.60\textwidth}
Find the maximum in \{-4,-9,-2\}. State the policy for an empty array.
\end{column}\begin{column}{0.35\textwidth}\small
Attempt the solution before continuing.

Use the definitions and examples from this lecture.
\end{column}\end{columns}
\note{Give students 8 minutes to attempt the problem. The following slides provide a complete solution. Original solution guidance: Initialise max from the first element, then compare remaining values. Result: -2. Reject empty input or return an explicitly designed optional result.\par Syllabus reading: Deitel Ch 7. CLO-2.}
\end{frame}

\begin{frame}[fragile,t]{Solution strategy}
\begin{columns}[T,onlytextwidth]\begin{column}{0.60\textwidth}
\begin{itemize}
\item Require a nonempty array before accessing values[0].
\item Initialise the maximum from the first element.
\item Visit each remaining element and replace the candidate only when the new value is greater.
\end{itemize}
\end{column}\begin{column}{0.35\textwidth}\small
The next program uses fixed inputs so its result can be reproduced.
\end{column}\end{columns}
\note{Discuss why each step is needed. State the input assumptions before executing the program.\par Syllabus reading: Deitel Ch 7. CLO-2.}
\end{frame}

\begin{frame}[fragile,t]{Complete solution (1/2)}
\begin{lstlisting}
public class Solution20 {
  public static void main(String[] args) throws Exception {
    int[] values = {-4, -9, -2};
    System.out.println(maximum(values));
  }
  static int maximum(int[] values) {
    if (values == null || values.length == 0) {
      throw new IllegalArgumentException("nonempty array required");
\end{lstlisting}
\note{Complete runnable file: Java\_Examples/Solution20.java. Inputs are fixed in the program.  \par Syllabus reading: Deitel Ch 7. CLO-2.}
\end{frame}

\begin{frame}[fragile,t]{Complete solution (2/2)}
\begin{lstlisting}
    }
    int maximum = values[0];
    for (int i = 1; i < values.length; i++) {
      if (values[i] > maximum) maximum = values[i];
    }
    return maximum;
  }
}
\end{lstlisting}
\note{Complete runnable file: Java\_Examples/Solution20.java. Inputs are fixed in the program.  \par Syllabus reading: Deitel Ch 7. CLO-2.}
\end{frame}

\begin{frame}[fragile,t]{Execution trace}
\small\renewcommand{\arraystretch}{1.5}
\rowcolors{2}{pale}{white}
\begin{tabularx}{\textwidth}{>{\raggedright\arraybackslash}X>{\raggedright\arraybackslash}X>{\raggedright\arraybackslash}X>{\raggedright\arraybackslash}X}
\rowcolor{navy}
\color{white}\bfseries Step & \color{white}\bfseries Compared value & \color{white}\bfseries Previous maximum & \color{white}\bfseries New maximum\\
Initialisation & values[0] = -4 & None & -4\\
i = 1 & -9 & -4 & -4\\
i = 2 & -2 & -4 & -2\\
\end{tabularx}
\vspace{3mm}\footnotesize Follow the changing state in execution order.
\note{Manually derived trace for the complete solution. Require a nonempty array before accessing values[0]. Initialise the maximum from the first element. Visit each remaining element and replace the candidate only when the new value is greater.\par Syllabus reading: Deitel Ch 7. CLO-2.}
\end{frame}

\begin{frame}[fragile,t]{Solution result}
\begin{columns}[T,onlytextwidth]\begin{column}{0.60\textwidth}
-2
\end{column}\begin{column}{0.35\textwidth}\small
Why this result follows

Visit each remaining element and replace the candidate only when the new value is greater.
\end{column}\end{columns}
\note{Expected result: -2\par Syllabus reading: Deitel Ch 7. CLO-2.}
\end{frame}

\begin{frame}[fragile,t]{A common incorrect approach}
for(int i=0; i\textless{}=a.length; i++) \{\par   System.out.println(a[i]);\par \}
\vspace{1mm}\begin{block}{Example}
\begin{lstlisting}
Use i<a.length. The last valid index is length-1.
\end{lstlisting}
\end{block}
\note{Ask students to predict the failure before discussing the explanation. The right side gives the correction.\par Syllabus reading: Deitel Ch 7. CLO-2.}
\end{frame}

\begin{frame}[fragile,t]{Boundary and failure checks}
\begin{columns}[T,onlytextwidth]\begin{column}{0.60\textwidth}
Check the stated input assumptions.

Read five temperatures into an array. Print the mean and number above the mean.
\end{column}\begin{column}{0.35\textwidth}\small
Initialise max from the first element, then compare remaining values. Result: -2. Reject empty input or return an explicitly designed optional result.
\end{column}\end{columns}
\note{Use the specification to derive each expected result. Exercise solution: Initialise max from the first element, then compare remaining values. Result: -2. Reject empty input or return an explicitly designed optional result.\par Syllabus reading: Deitel Ch 7. CLO-2.}
\end{frame}

\begin{frame}[fragile,t]{Check your understanding}
\begin{columns}[T,onlytextwidth]\begin{column}{0.60\textwidth}
What values does new int[3] contain? Does changing the enhanced-for variable update the array?
\end{column}\begin{column}{0.35\textwidth}\small
Write a prediction and explain the rule that supports it.
\end{column}\end{columns}
\note{Answer: Three zeros. No, for a primitive component it only changes the loop variable.\par Syllabus reading: Deitel Ch 7. CLO-2.}
\end{frame}

\begin{frame}[fragile,t]{Answer and explanation}
\begin{columns}[T,onlytextwidth]\begin{column}{0.60\textwidth}
Three zeros. No, for a primitive component it only changes the loop variable.
\end{column}\begin{column}{0.35\textwidth}\small
Use i\textless{}a.length. The last valid index is length-1.
\end{column}\end{columns}
\note{Reconnect the answer to the relevant definition rather than treating it as a fact to memorise.\par Syllabus reading: Deitel Ch 7. CLO-2.}
\end{frame}


\begin{frame}[fragile,t]{Compare cases and predict outcomes}
\small\renewcommand{\arraystretch}{1.5}\rowcolors{2}{pale}{white}
\begin{tabularx}{\textwidth}{>{\raggedright\arraybackslash}X>{\raggedright\arraybackslash}X>{\raggedright\arraybackslash}X}\rowcolor{navy}
\color{white}\bfseries Index & \color{white}\bfseries Element & \color{white}\bfseries Running sum after visit\\
0 & 12 & 12\\
1 & 7 & 19\\
2, then 3 & 19, then 4 & 38, then 42\\
\end{tabularx}\par\vspace{3mm}\begin{exampleblock}{Discuss}Explain each expected result before running the example. Which case would reveal a wrong assumption?\end{exampleblock}
\note{Read the first two columns and invite predictions before discussing the outcome.}
\end{frame}

\begin{frame}[fragile,t]{Apply the idea}
\begin{alertblock}{Independent practice}Write a loop to count the values above 10 in \{12,7,19,4\}. Give the count and identify the indexes that qualify.\end{alertblock}\vspace{3mm}\begin{enumerate}\item Work through the input by hand.\item Write or correct the program.\item Justify the expected result.\end{enumerate}\note{Allow 3–5 minutes, then compare answers in pairs. Write a loop to count the values above 10 in \{12,7,19,4\}. Give the count and identify the indexes that qualify.}
\end{frame}

\begin{frame}[fragile,t]{Solution and reasoning}
\begin{lstlisting}
int[] values = {12, 7, 19, 4};
int count = 0;
for (int value : values) {
  if (value > 10) count++;
}
System.out.println(count);
\end{lstlisting}\vspace{1mm}\small\begin{itemize}
\item Values 12 and 19 qualify, at indexes 0 and 2.
\item The output is 2.
\item An enhanced for loop is sufficient because no element needs modification.
\end{itemize}\note{Answer: Values 12 and 19 qualify, at indexes 0 and 2. The output is 2. An enhanced for loop is sufficient because no element needs modification.}
\end{frame}

\begin{frame}[fragile,t]{Linear search and the first match}
\begin{itemize}
\item Compare the target with successive elements until the first match.
\item Return its index, not the element value. Use {-}1 when no match exists.
\item Linear search works without sorting; binary search has a sorted{-}input precondition.
\end{itemize}
\vspace{3mm}\footnotesize\textcolor{accent}{Source: supplied ISB campus slides. See speaker notes and ISB\_Source\_Map.md.}
\note{Adapted from the supplied ISB campus folder: Arrays1D(new).pptx, slides 30, 31, 32; Logic Building, Dry Run Techniques, and Flagged While Loops.pptx, slides 13. Adapted explanation and runnable implementation; sample inputs and traces added for this course.}
\end{frame}

\begin{frame}[fragile,t]{Worked problem: Linear search and the first match}
\begin{block}{Task}Search \{5,8,12,3,9\} for 12 and for 7, using the source search contract.\end{block}
\small Input: Values are fixed in the program for a reproducible demonstration.\par\vspace{3mm}\begin{exampleblock}{Before running}Predict the output and identify the variables that change.\end{exampleblock}
\note{Adapted from the supplied ISB campus folder: Arrays1D(new).pptx, slides 30, 31, 32; Logic Building, Dry Run Techniques, and Flagged While Loops.pptx, slides 13. Adapted explanation and runnable implementation; sample inputs and traces added for this course.}
\end{frame}

\begin{frame}[fragile,t]{Complete ISB example (1/1)}
\begin{lstlisting}
public class IsbLecture20 {
  public static void main(String[] args) throws Exception {
    int[] values = {5, 8, 12, 3, 9};
    System.out.println(linearSearch(values, 12));
    System.out.println(linearSearch(values, 7));
  }
  static int linearSearch(int[] values, int key) {
    for (int i = 0; i < values.length; i++) {
      if (values[i] == key) return i;
    }
    return -1;
  }
}
\end{lstlisting}
\note{Adapted from the supplied ISB campus folder: Arrays1D(new).pptx, slides 30, 31, 32; Logic Building, Dry Run Techniques, and Flagged While Loops.pptx, slides 13. Adapted explanation and runnable implementation; sample inputs and traces added for this course. Complete file: ISB\_Java\_Examples/IsbLecture20.java.}
\end{frame}

\begin{frame}[fragile,t]{Worked trace and expected result}
\small\renewcommand{\arraystretch}{1.4}\rowcolors{2}{pale}{white}
\begin{tabularx}{\textwidth}{>{\raggedright\arraybackslash}X>{\raggedright\arraybackslash}X>{\raggedright\arraybackslash}X}\rowcolor{navy}
\color{white}\bfseries Index visited & \color{white}\bfseries Value & \color{white}\bfseries Searching for 12\\
0 & 5 & Continue\\
1 & 8 & Continue\\
2 & 12 & Return 2\\
\end{tabularx}\par\vspace{2mm}
\begin{block}{Expected console output}\ttfamily\footnotesize 2\par {-}1\end{block}
\note{Adapted from the supplied ISB campus folder: Arrays1D(new).pptx, slides 30, 31, 32; Logic Building, Dry Run Techniques, and Flagged While Loops.pptx, slides 13. Adapted explanation and runnable implementation; sample inputs and traces added for this course. Trace and expected values independently worked for this edition.}
\end{frame}

\begin{frame}[fragile,t]{Extend the ISB example}
\begin{alertblock}{Practice}Predict searching \{12,5,12\} for 12 and searching an empty array.\end{alertblock}\vspace{3mm}\begin{exampleblock}{Explain your reasoning}The first returns 0; the empty array returns {-}1 without entering the loop.\end{exampleblock}
\note{Adapted from the supplied ISB campus folder: Arrays1D(new).pptx, slides 30, 31, 32; Logic Building, Dry Run Techniques, and Flagged While Loops.pptx, slides 13. Adapted explanation and runnable implementation; sample inputs and traces added for this course. Answer: The first returns 0; the empty array returns {-}1 without entering the loop.}
\end{frame}
\begin{frame}[fragile,t]{Lecture summary}
\begin{columns}[T,onlytextwidth]\begin{column}{0.60\textwidth}
\begin{itemize}
\item arrays with fixed or runtime sizes
\item Traverse elements safely
\item an aggregate and explain bounds errors
\end{itemize}
\end{column}\begin{column}{0.35\textwidth}\small
\begin{itemize}
\item Require a nonempty array before accessing values[0].
\item Initialise the maximum from the first element.
\item Visit each remaining element and replace the candidate only when the new value is greater.
\end{itemize}
\end{column}\end{columns}
\note{Ask students to give one concrete example for the concepts listed. Use the worked solution to connect the topics.\par Syllabus reading: Deitel Ch 7. CLO-2.}
\end{frame}

\begin{frame}[fragile,t]{After-class practice}
\begin{columns}[T,onlytextwidth]\begin{column}{0.60\textwidth}
Read five temperatures into an array. Print the mean and number above the mean.
\end{column}\begin{column}{0.35\textwidth}\small
Syllabus reading\par Deitel Ch 7

Next lecture: Arrays and method arguments
\end{column}\end{columns}
\note{Reading labels follow the supplied outline. Check topic headings against the available textbook edition. Require a program and expected results for the practice task.\par Syllabus reading: Deitel Ch 7. CLO-2.}
\end{frame}
\end{document}
